\ticket{25}{Извлечение информации из текстов: особенности задачи, виды извлекаемых данных, применяемые подходы. Задача извлечения мнений и анализа тональности текстов}

\begin{examplan}
    \item Определить извлечение информации из текстов.
    \item Перечислить виды извлекаемых данных.
    \item Сравнить подходы к извлечению информации.
    \item Описать извлечение мнений и анализ тональности.
\end{examplan}

\subsection*{Короткая версия для письменного ответа}

\textbf{Извлечение информации} --- автоматическое преобразование неструктурированного текста в структурированные факты: сущности, отношения, события, шаблоны или мнения. В отличие от информационного поиска, IE извлекает не документы, а конкретные данные из текста.

\subsection*{Особенности задачи}

\begin{itemize}
    \item исходные тексты не имеют строгой структуры;
    \item заранее задаются типы нужной информации;
    \item важны лексическая вариативность, синтаксис, контекст и неоднозначность;
    \item системы часто зависят от предметной области;
    \item качество оценивается через precision, recall и F1.
\end{itemize}

\subsection*{Виды извлекаемых данных}

\begin{description}
    \item[Именованные сущности] люди, организации, локации, даты, денежные суммы. Задача называется NER.
    \item[Отношения] связи между сущностями, например works-for или located-in.
    \item[События] тип события, триггер и аргументы с ролями.
    \item[Шаблоны] заполнение заранее заданной структуры, например карточки происшествия.
    \item[Мнения] объект, аспект, держатель мнения, полярность и интенсивность.
\end{description}

\subsection*{Подходы к IE}

\textbf{Правиловый подход} использует регулярные выражения, газеттиры, лексико-синтаксические шаблоны и грамматики. Он интерпретируем и точен для узких задач, но плохо переносится на новые домены.

\textbf{ML-подход} обучается на размеченных корпусах. NER часто формулируется как разметка последовательности с BIO-метками. Используются HMM, CRF, BiLSTM-CRF, Transformer-модели. Relation Extraction часто решается как классификация пары сущностей в контексте.

\textbf{Гибридные подходы} совмещают правила и модели: например, газеттиры используются как признаки, а ML-модель фильтрует кандидатов.

\textbf{Distant supervision} автоматически создает шумную разметку из баз знаний: если две сущности связаны отношением в базе, предложения с этой парой считаются кандидатами для обучения.

\textbf{OpenIE} извлекает открытые тройки вида аргумент--отношение--аргумент без заранее фиксированного набора отношений.

\subsection*{Анализ тональности}

\textbf{Анализ тональности} выявляет субъективную оценку в тексте: позитивную, негативную, нейтральную или более детальную шкалу.

Уровни:
\begin{itemize}
    \item уровень документа;
    \item уровень предложения;
    \item аспектный уровень.
\end{itemize}

\textbf{Аспектный анализ} отвечает на вопросы: кто выражает мнение, о чем, по какому аспекту и с какой тональностью. Пример: "Камера отличная, но батарея слабая" --- камера позитивно, батарея негативно.

Подходы:
\begin{itemize}
    \item словари тональности: подсчет позитивных и негативных слов с учетом отрицаний и усилителей;
    \item машинное обучение: BoW, n-граммы, эмбеддинги, Naive Bayes, логистическая регрессия, SVM, CNN/RNN/Transformer;
    \item нейросетевые модели для аспектов, контекста, сарказма и сложных формулировок.
\end{itemize}

Сложности: сарказм, ирония, доменная зависимость лексики, отрицания, смешанные оценки и неявные мнения.

\begin{important}
    \item IE извлекает структурированные факты из неструктурированного текста.
    \item Основные виды: сущности, отношения, события, шаблоны и мнения.
    \item NER часто решается как sequence labeling.
    \item Анализ тональности может быть документным, предложенческим или аспектным.
    \item Правила интерпретируемы, ML лучше масштабируется при наличии данных.
\end{important}
